% 原书物理页=印刷页44--51；章2标题起，2.2标题前止.

\chapter{几何与算术理论}\label{chap:modular-2}

\section{模形式与Galois表示的联系}\label{sec:modular-2-1}

我们的目标是尽可能解释模形式与Galois表示之间的一些联系.数学家发现，这两个概念都会产生$L$函数，而这些$L$函数的性质，无论已获证明还是尚属猜想，都十分相似；二者的联系正是由此显现的.



\subsection{Galois表示}\label{sec:modular-2-1-1}

$\mathbb Q$的绝对Galois群是数论中的基本研究对象.与研究任何群一样，研究它的线性表示是认识它的一个好办法.因此，我们考察连续表示
\begin{align*}
\rho:\operatorname{Gal}(\overline{\mathbb Q}/\mathbb Q)\longrightarrow\operatorname{GL}_n(C),
\end{align*}
其中$C$表示特征为$0$的局部域（也可以是这种域的代数闭包）.这里，$\operatorname{GL}_n(C)$带有使它成为$\operatorname M_n(C)$中开集的拓扑，而$\operatorname M_n(C)$是局部域$C$上的有限维向量空间.群$\operatorname{Gal}(\overline{\mathbb Q}/\mathbb Q)$则带有如下拓扑：$\operatorname{Gal}(\overline{\mathbb Q}/K)$（$K$遍历$\mathbb Q$的有限扩张）构成单位元的开邻域基.将该群写成
\begin{align*}
\operatorname{Gal}(\overline{\mathbb Q}/\mathbb Q)
=\left\{(\gamma_K)_K\in\prod_{[K:\mathbb Q]<\infty}\operatorname{Gal}(K/\mathbb Q):
 (\gamma_K)|_{K'}=\gamma_{K'}\text{，若 }K'\subset K\right\},
\end{align*}
这也是乘积拓扑.由此可知，$\operatorname{Gal}(\overline{\mathbb Q}/\mathbb Q)$是``投有限''的（有限群的投射极限），因而是紧的.\footnote{译注：上式沿用原文；这一逆极限描述应让$K$遍历$\mathbb Q$的有限Galois扩张.任意有限扩张未必在绝对Galois群作用下稳定.}于是有如下两种情形：

\begin{enumerate}[label=\roman*)]

\item 若$C=\mathbb C$（或$\mathbb R$），则$\rho$连续，当且仅当它经某个有限Galois群$\operatorname{Gal}(K/\mathbb Q)$分解.这些表示称为\newterm{modular-artin-representation}{Artin表示}{Artin representation}.

\item 若$C$是域$\mathbb Q_\ell$（$\mathbb Q$关于$\ell$进绝对值的完备化）的扩张，则$\rho$的连续性等价于下列性质：

\begin{itemize}

\item 必要时在$\operatorname{GL}_2(C)$中作共轭后，$\rho$取值于$\operatorname{GL}_2(\mathcal O_C)$.

\item 对每个$n$，将表示$\rho$模$\ell^n$约化后，它经$\mathbb Q$的某个有限扩张的Galois群分解.

\end{itemize}

这些表示通常简称为\newterm{modular-l-adic-representation}{$\ell$进表示}{$\ell$-adic representation}.\footnote{译注：原文在一般的$n$维表示之后使用$\operatorname{GL}_2$，此处保留.关于取值于整数环的断言，通常在$C/\mathbb Q_\ell$为有限扩张等适当条件下表述.}

\end{enumerate}



\begin{example*}[$n=1$：分圆特征]

设$N$为整数，$\mu_N$为$\overline{\mathbb Q}$中$N$次单位根组成的群，$\zeta_N\in\mu_N$为本原根.回忆，映射$\sigma\mapsto n_\sigma$，其中$\sigma(\zeta_N)=\zeta_N^{n_\sigma}$，给出同构
\begin{align*}
\varepsilon_N:\operatorname{Gal}(\mathbb Q(\mu_N)/\mathbb Q)\xrightarrow{\sim}(\mathbb Z/N\mathbb Z)^\times.
\end{align*}
因此，每个模$N$的Dirichlet特征$\chi$都给出连续特征
\begin{align*}
\chi\circ\varepsilon_N:\operatorname{Gal}(\overline{\mathbb Q}/\mathbb Q)
\longrightarrow\operatorname{Gal}(\mathbb Q(\mu_N)/\mathbb Q)
\longrightarrow\operatorname{GL}_1(\mathbb C).
\end{align*}
另一方面，$\varepsilon_{\ell^n}$（$n\in\mathbb N$）这一族同构在极限中给出同构
\begin{align*}
\varepsilon_{\ell^\infty}:\operatorname{Gal}(\mathbb Q(\mu_{\ell^\infty})/\mathbb Q)
\xrightarrow{\sim}\varprojlim_n(\mathbb Z/\ell^n\mathbb Z)^\times=\mathbb Z_\ell^\times,
\end{align*}
因而可以将它看成连续特征
\begin{align*}
\varepsilon_{\ell^\infty}:\operatorname{Gal}(\overline{\mathbb Q}/\mathbb Q)
\longrightarrow\operatorname{Gal}(\mathbb Q(\mu_{\ell^\infty})/\mathbb Q)
\longrightarrow\operatorname{GL}_1(\mathbb Q_\ell).
\end{align*}
Kronecker--Weber定理告诉我们，这样便构造出了$\operatorname{Gal}(\overline{\mathbb Q}/\mathbb Q)$的所有一维连续表示.

\end{example*}



\begin{example*}[$n=2$：椭圆曲线]

设$E$为$\mathbb Q$上的椭圆曲线（稍后将回顾一些有关知识）.已知$N$挠点子群$E(\overline{\mathbb Q})[N]$同构于$(\mathbb Z/N\mathbb Z)^2$.它显然还带有与$E$中的乘法相容的$\operatorname{Gal}(\overline{\mathbb Q}/\mathbb Q)$作用，而且这些$N$挠点定义在$\mathbb Q$的某个有限扩张上.因此得到连续表示
\begin{align*}
\rho_{E,N}:\operatorname{Gal}(\overline{\mathbb Q}/\mathbb Q)\longrightarrow\operatorname{GL}_2(\mathbb Z/N\mathbb Z).
\end{align*}
Weil配对告诉我们，$\rho_{E,N}$的行列式由$\varepsilon_N$给出.

考虑$E(\overline{\mathbb Q})[\ell^n]$组成的投射系，转移映射取为乘以$\ell$.取极限便得到连续表示
\begin{align*}
\rho_{E,\ell^\infty}:\operatorname{Gal}(\overline{\mathbb Q}/\mathbb Q)
\longrightarrow\varprojlim_n\operatorname{GL}_2(\mathbb Z/\ell^n\mathbb Z)
=\operatorname{GL}_2(\mathbb Z_\ell)\subset\operatorname{GL}_2(\mathbb Q_\ell).
\end{align*}
\end{example*}



\begin{example*}[$\ell$进上同调]

更一般地，对定义在$\mathbb Q$上的代数簇$X$，Grothendieck定义了``平展上同调''群$H^i_{\text{ét}}(X_{\overline{\mathbb Q}},\mathbb Z/N\mathbb Z)$，Galois群$\operatorname{Gal}(\overline{\mathbb Q}/\mathbb Q)$作用在这些群上.对$H^i_{\text{ét}}(X_{\overline{\mathbb Q}},\mathbb Z/\ell^n\mathbb Z)$关于$n$取投射极限，再将$\ell$逆化，便得到$\mathbb Q_\ell$上的有限维向量空间$H^i_{\text{ét}}(X_{\overline{\mathbb Q}},\mathbb Q_\ell)$，其上带有$\operatorname{Gal}(\overline{\mathbb Q}/\mathbb Q)$的连续线性作用.这样的$\ell$进表示称为``源于几何的''表示.这里使用$\ell$进系数至关重要：Serre指出，不存在以复数为系数、对Galois作用等变的上同调理论.

\end{example*}



\paragraph{局部表示.}

设$p$为素数.固定$\mathbb Q_p$的代数闭包$\overline{\mathbb Q}_p$.每个嵌入$\iota:\overline{\mathbb Q}\hookrightarrow\overline{\mathbb Q}_p$，通过限制自同构，诱导出嵌入
\begin{align*}
\iota_*:G_{\mathbb Q_p}:=\operatorname{Gal}(\overline{\mathbb Q}_p/\mathbb Q_p)
\hookrightarrow\operatorname{Gal}(\overline{\mathbb Q}/\mathbb Q).
\end{align*}
由于任何其他的$\iota'$都可写成$\iota\circ\gamma$，其中$\gamma\in\operatorname{Gal}(\overline{\mathbb Q}/\mathbb Q)$，故嵌入$\iota'_*$与$\iota_*$在$\gamma$下共轭.因此，$\rho$在$G_{\mathbb Q_p}$上的两个限制$\rho\circ\iota_*$和$\rho\circ\iota'_*$在$\rho(\gamma)$下共轭，因而等价.所得$\operatorname{Gal}(\overline{\mathbb Q}_p/\mathbb Q_p)$表示的等价类，简记为$\rho|_{G_{\mathbb Q_p}}$.



\paragraph{分歧.}

记$|\cdot|_p$为$p$进绝对值在$\overline{\mathbb Q}_p$上的唯一延拓.$G_{\mathbb Q_p}$在$\overline{\mathbb Q}_p$上的作用保持绝对值$|\cdot|_p$，特别地，保持整数环
\begin{align*}
\overline{\mathbb Z}_p&=\{x\in\overline{\mathbb Q}_p:|x|_p\leqslant1\},&
\overline{\mathfrak m}_p&:=\{x\in\overline{\mathbb Q}_p:|x|_p<1\}
\end{align*}
及其极大理想$\overline{\mathfrak m}_p$.于是，$G_{\mathbb Q_p}$以域自同构作用在商$\overline{\mathbb Z}_p/\overline{\mathfrak m}_p=\overline{\mathbb F}_p$上，从而得到同态
\begin{align*}
G_{\mathbb Q_p}\longrightarrow G_{\mathbb F_p}=\operatorname{Gal}(\overline{\mathbb F}_p/\mathbb F_p).
\end{align*}
将其核记为$I_{\mathbb Q_p}$，称为$G_{\mathbb Q_p}$的\newterm{modular-inertia-subgroup}{惯性子群}{inertia subgroup}.若$\rho|_{I_{\mathbb Q_p}}$平凡，即$I_{\mathbb Q_p}$在表示$\rho$的空间$V=C^n$上以恒等变换作用，就称表示$\rho$在$p$处非分歧（也称$\rho|_{G_{\mathbb Q_p}}$非分歧）.

\begin{exercise*}

$p$在$\rho$中非分歧，当且仅当$\rho$经$\operatorname{Gal}(\overline{\mathbb Q}^{(p')}/\mathbb Q)$分解，其中$\overline{\mathbb Q}^{(p')}$是$\overline{\mathbb Q}$中使$p$非分歧的最大子域.[提示：这等价于证明，$\overline{\mathbb Q}^{(p')}$是由所有$\iota_*(I_{\mathbb Q_p})$生成的子群的不动子域，其中$\iota$遍历嵌入$\overline{\mathbb Q}\hookrightarrow\overline{\mathbb Q}_p$.]

\end{exercise*}

我们只考虑仅有有限多个分歧素数的表示.Artin表示自动满足这一条件：例如，若$\rho$经$\operatorname{Gal}(K/\mathbb Q)$分解，则每个在$K$中非分歧的$p$对$\rho$也非分歧.相反，$\ell$进表示并不自动满足这一条件.对整数$N$，若$\rho$中分歧的素数只能是$N$的素因子，就称$\rho$在$N$之外非分歧.这等价于要求$\rho$经$\operatorname{Gal}(\overline{\mathbb Q}^{(N)}/\mathbb Q)$分解，其中$\overline{\mathbb Q}^{(N)}$表示$\overline{\mathbb Q}$中仅可能在$N$的素因子处分歧的最大子域.



\paragraph{Frobenius特征多项式.}

已知同态$G_{\mathbb Q_p}\longrightarrow G_{\mathbb F_p}$是满射.任取$G_{\mathbb Q_p}$中映到Frobenius$(x\mapsto x^p)\in G_{\mathbb F_p}$的元素$\sigma_p$.那么$\sigma_p$正规化$I_{\mathbb Q_p}$，从而保持表示$\rho$的空间$V=C^n$中的$I_{\mathbb Q_p}$不变子空间$V^{I_{\mathbb Q_p}}$.于是定义
\begin{align*}
\Phi_{p,\rho}(X):=\det\left((\mathrm{id}-X\rho(\sigma_p))\big|_{V^{I_{\mathbb Q_p}}}\right)\in C[X],
\end{align*}
这是$\rho(\sigma_p)$作用在该$I_{\mathbb Q_p}$不变子空间上的特征多项式（或者说，几乎是特征多项式）.它不依赖$\sigma_p$的选取，因为其他任何选择都属于$\sigma_p I_{\mathbb Q_p}$.由定义可见，$\Phi_p(X)$的次数至多为$n$，而$p$分歧当且仅当$\Phi_p(X)$的次数小于$n$.



\subsection{Artin表示与权为1的形式}\label{sec:modular-2-1-2}

从Artin表示出发，更容易看出模形式与Galois表示之间存在对应的可能性.这时，多项式$\Phi_p(X)$的复系数有与$p$无关的界.事实上，$\rho$经一个有限群分解，所以$\rho(\sigma_p)$的特征值都是单位根，从而
\begin{align*}
\left|\operatorname{tr}\left(\bigwedge^i\rho(\sigma_p)\right)\right|\leqslant\binom ni.
\end{align*}
因此Euler乘积
\begin{align*}
L(s,\rho)=\prod_p\Phi_p(p^{-s})^{-1}
\end{align*}
在$\Re(s)>1$时收敛.

\begin{exercise*}

设$n=1$.若$\rho$是平凡表示，则$L(s,\rho)=\zeta(s)$.更一般地，若$\rho=\chi\circ\varepsilon_N$，其中$\chi$是模$N$的本原Dirichlet特征，则
\begin{align*}
L(s,\rho)=L(s,\chi):=\sum_{n=1}^\infty\frac{\chi(n)}{n^s}.
\end{align*}
\end{exercise*}

\begin{example*}

设$K$为$\mathbb Q$的有限扩张，$X$为$K$到$\overline{\mathbb Q}$的嵌入组成的有限集.群$\operatorname{Gal}(\overline{\mathbb Q}/\mathbb Q)$作用在$X$上，因而线性地作用在向量空间$\mathbb C^X$上.由此得到一个维数为$n=[K:\mathbb Q]$的Artin表示$\rho$.可以证明（这是一个很好的习题），
\begin{align*}
L(s,\rho)=\zeta_K(s)
=\prod_{\mathfrak p\in\operatorname{Max}(\mathcal O_K)}\left(1-(N\mathfrak p)^{-s}\right)^{-1}
=\sum_{\mathfrak a}|N\mathfrak a|^{-s}
\end{align*}
（即数域$K$的$\zeta$函数）.这里$\mathfrak a$遍历$\mathcal O_K$的非零理想.

\end{example*}

\begin{remark*}

若$\rho=\rho_1\oplus\rho_2$，则$L(s,\rho)=L(s,\rho_1)L(s,\rho_2)$.

\end{remark*}

Chebotarev定理告诉我们，有限Galois群$\operatorname{Gal}(K/\mathbb Q)$的每个元素都是某个在$K$中非分歧的素数$p$处的Frobenius代换\footnote{即形如$\iota_*(\sigma_p)|_K$的元素，其中$\iota$为一个嵌入，$\sigma_p$为上文所说的Frobenius提升.}（甚至对应于无穷多个这样的$p$）.有限群的复表示由其特征标确定，所以Artin表示$\rho$由$\rho$非分歧的素数$p$处的$\operatorname{tr}(\rho(\sigma_p))$唯一确定.而在展开式$L(s,\rho)=\sum_{n=1}^\infty b_n n^{-s}$中，可以恢复出$\operatorname{tr}(\rho(\sigma_p))=b_p$，所以$\rho$更由其函数$L(s,\rho)$确定.

Artin证明了，$L(s,\rho)$可以亚纯延拓到$s\in\mathbb C$，并满足某种函数方程，其类型与$\zeta$、$L(s,\chi)$或$\zeta_K$所满足的函数方程相同.随后他猜想，这一延拓$L(s,\rho)$全纯，当且仅当$\rho$在$\operatorname{Gal}(\overline{\mathbb Q}/\mathbb Q)$作用下没有非零不变向量.这个猜想仍是一个未解决的问题.

当$\rho$为二维时，$L(s,\rho)$的函数方程与权为$1$的本原形式的$L$函数$L(s,f)$所满足的函数方程相似.此外，Hecke已经证明了$L(s,f)$的全纯性.下面这个引人瞩目的定理曾由Langlands猜想提出，他从中看到了Artin猜想的深刻缘由.

\begin{theorem*}[Deligne--Serre]

设$f\in\mathcal S_1(N,\chi)$为本原形式.存在唯一的二维Artin表示$\rho_f$，使得$L(s,f)=L(s,\rho_f)$.

\end{theorem*}

\begin{remark*}

等式$L(s,f)=L(s,\rho_f)$等价于所有局部因子相等：
\begin{align*}
\forall p,\qquad1-a_p(f)p^{-s}+\chi(p)p^{-2s}=\Phi_p(p^{-s}),
\end{align*}
也就等价于
\begin{align*}
&\forall p\nmid N,\quad a_p(f)=\operatorname{tr}(\rho(\sigma_p)),\quad
\chi(p)=\det(\rho(\sigma_p)),\\
&\forall p\mid N,\quad a_p(f)=\operatorname{tr}\left(\rho(\sigma_p)\big|_{V^{I_p}}\right).
\end{align*}
特别地，$\rho_f$在$N$之外非分歧，而在每个整除$N$的素数处分歧.实际上，如上所见，表示$\rho_f$完全由下列等式确定：
\begin{align*}
\operatorname{tr}(\rho(\sigma_p))=a_p(f)\qquad(p\nmid N).
\end{align*}
Deligne与Serre在证明中先构造$\rho_f$，使$p\nmid N$处的局部因子相等，再利用两个$L$函数各自的函数方程，证明$p\mid N$处局部因子的相等性.

\end{remark*}

\begin{remark*}

因为$L(s,f)$不是两个Dirichlet $L$函数的乘积，所以表示$\rho_f$不可约.另一方面，其行列式显然为$\det(\rho_f)=\chi\circ\varepsilon_N$.特别地，复共轭\footnote{为简便起见，可以假设$\overline{\mathbb Q}$已经嵌入$\mathbb C$.否则，复共轭只在相差共轭的意义下确定；它对于阿基米德绝对值起着Frobenius代换的作用.}$\sigma_\infty$满足
\begin{align*}
\det(\rho(\sigma_\infty))=\chi(-1)=-1.
\end{align*}
最后这个等式是$\mathcal S_1(N,\chi)\ne0$的必要条件.

\end{remark*}

值得注意的是，这个结果使两个领域都得到了新的结论：

\begin{itemize}

\item 一方面，它推出权为$1$的$f$满足Ramanujan--Petersson猜想所预言的$|a_p(f)|\leqslant2$.

\item 另一方面，它对表示$\rho_f$推出Artin猜想.事实上，Langlands赋予Artin猜想的现代表述是：每个二维不可约表示都形如$\rho_f$.\footnote{译注：原文此处未重复``Artin''及``奇''的限制；由上文$\det(\rho_f(\sigma_\infty))=-1$可见，这种对应至少需要奇性条件.此处保留原文陈述.}注意，对二维可约表示，即$\rho=(\chi\circ\varepsilon_N)\oplus(\chi'\circ\varepsilon_{N'})$，已知如何构造Eisenstein级数，使$L(s,f)=L(s,\chi)L(s,\chi')$.

\end{itemize}

这里不讲解Deligne与Serre的证明，但要指出：虽然定理涉及的对象``较简单''（Artin表示），其证明却用到了与权$k>1$的形式相联系的Galois表示的构造；如下所见，后者是$\ell$进表示.



\subsection{$\ell$进表示与权为$k\geqslant2$的形式}\label{sec:modular-2-1-3}

\begin{theorem*}[Deligne、Shimura]

设$f\in\mathcal S_k(N,\chi)$为权$k$、级$N$、特征$\chi$的本原形式.

\begin{enumerate}[label=\roman*)]

\item 由特征值$\lambda(n)$（$n\in\mathbb N$）在$\mathbb Q$上生成的域$K_f$是$\mathbb Q$的有限次扩张，且各$\lambda(n)$都是代数整数.

\item 设$\mathfrak l$为$\mathcal O_{K_f}$中位于素数$\ell$之上的极大理想.记$K_{f,\mathfrak l}$为$K_f$关于$\mathfrak l$所对应绝对值的完备化.因此，它是$\mathbb Q_\ell$的有限扩张.那么，存在唯一的（不计共轭）连续表示
\begin{align*}
\rho_{f,\mathfrak l}:\operatorname{Gal}(\overline{\mathbb Q}/\mathbb Q)\longrightarrow\operatorname{GL}_2(K_{f,\mathfrak l}),
\end{align*}
使得对每个$p\nmid\ell N$，都有
\begin{align*}
(*)_p\qquad\Phi_{p,\rho_{f,\mathfrak l}}(X)=1-\lambda(p)X+\chi(p)p^{k-1}X^2.
\end{align*}
此外，$\rho_{f,\mathfrak l}$不可约，其行列式为特征$(\chi\circ\varepsilon_N)\varepsilon_{\ell^\infty}^{k-1}$，而等式$(*)_p$对每个$p\ne\ell$都成立.

\end{enumerate}

\end{theorem*}

\begin{remark*}[关于唯一性]

若表示$\rho_{f,\mathfrak l}$存在，则它显然在$p\nmid\ell N$处非分歧（而在$p\mid\ell N$处分歧）.\footnote{译注：括号内为原文断言.非分歧素数处的等式本身不能推出每个$p\mid\ell N$都分歧，尤其$p=\ell$处需另作分析.}与Artin表示的情形相同，Chebotarev定理保证，$\rho_{f,\mathfrak l}$的半单化由各个$p\nmid\ell N$处的等式$(*)_p$确定，甚至仅由$\operatorname{tr}(\rho(\sigma_p))=\lambda(p)$（$p\nmid\ell N$）这些等式确定.而$\rho_{f,\mathfrak l}$不可约，故它完全由这些等式确定.

\end{remark*}

\begin{remark*}[关于行列式]

注意，由构造，$\chi\circ\varepsilon_N$确实取值于$K_{f,\mathfrak l}$，甚至取值于$K_f$.同样，$\varepsilon_{\ell^\infty}$取值于$\mathbb Z_\ell^\times\subset K_{f,\mathfrak l}^\times$.现在，假设$\rho_{f,\mathfrak l}$存在，根据$(*)_p$，$\det(\rho_{f,\mathfrak l})$与所声称的特征至少在所有$p\nmid\ell N$处的Frobenius代换上相等，因而由Chebotarev定理，它们在$\operatorname{Gal}(\overline{\mathbb Q}/\mathbb Q)$的每个元素上都相等.

\end{remark*}

\begin{remark*}[关于$p\mid N$且$p\ne\ell$时的等式$(*)_p$]

与权为$1$的情形相同，这些分歧素数处的等式可以由两侧$L$函数满足的函数方程推出.不过，先验并不清楚$\rho_{f,\mathfrak l}$是否具有一个合理的$L$函数（它并不是复表示！）.实际上，通过$\ell$进平展上同调来构造$\rho_{f,\mathfrak l}$，再结合这种上同调的一些非平凡性质，就可以为它配上这样的$L$函数.

\end{remark*}

\begin{remark*}[$p=\ell$时如何？]

$p=\ell$时的情形相当神秘，却对算术应用至关重要.为解释这一情形的特殊之处，注意当$p\nmid\ell N$时，可以把等式$(*)_p$解释为对限制$\rho|_{G_{\mathbb Q_p}}$的完整描述.但当$p\mid N$时，就出现了这样的奇特局面：我们知道$\rho|_{G_{\mathbb Q_p}}$已经确定（由Chebotarev定理），却不知道如何从$f$中``读出''它.即使$p\ne\ell$，等式$(*)_p$也不足以描述$\rho|_{G_{\mathbb Q_p}}$.然而在后一种情形，自Langlands以来，人们已经清楚如何从$f$对应的自守表示出发，完整地描述$\rho|_{G_{\mathbb Q_p}}$.这正是``局部Langlands对应''研究的内容.但是当$p=\ell$时，即使$f$对应的自守表示也不足以描述$\rho|_{G_{\mathbb Q_\ell}}$.近期的``$p$进Langlands纲领''旨在尝试理解这一情形，其主要工具是``$p$进Hodge理论''.

\end{remark*}

\begin{remark*}[关于定理的推论]

与权为$1$的情形相同，这个定理可以用来证明Ramanujan--Petersson猜想.不过，Galois一侧的估计要困难得多！这里同样要用到通过平展上同调所作的构造，以及Deligne证明的、关于有限域上代数簇点数大小的Weil猜想！

\end{remark*}

\begin{remark*}[关于逆命题]

与Artin表示的情形相同，这一定理有一个猜想性的逆命题.Langlands提出的第一个版本预言：每个``源于几何''的二维不可约表示（即出现为$\mathbb Q$上某个代数簇上同调的子商的表示）都来自模形式.这已经包含了由Shimura、Taniyama与Weil倡导的``椭圆曲线的模性猜想''；Wiles证明了它的一个特殊情形，从而完成了Fermat定理的证明.Fontaine与Mazur提出了一个更引人瞩目的版本，用一个与$\ell$进Hodge理论有关的条件取代``源于几何''的假设；详细陈述这个条件会使我们偏离主题太远.这个猜想（权$\geqslant2$的情形）最近取得了惊人的进展，只有少数情形仍未解决.但是……这还只是一个更一般猜想的$n=2$情形；后者联系着Galois表示与``自守表示''，距离理解它还很遥远……\footnote{译注：此段沿用原文的概述及其写作时的进展表述；二维几何表示的模性逆命题也需要适当的奇性与Hodge权等假设，不宜将此处概述视为无附加条件的精确陈述.}

\end{remark*}



\subsection{证明策略}\label{sec:modular-2-1-4}

第一个想法是找出一个有限生成自由Abel子群$H$，它生成$\mathbb C$向量空间$\mathcal S_k(N,\chi)$，并在Hecke环$\mathcal H_0(N)$的作用下稳定.Cayley--Hamilton定理于是告诉我们，Hecke算子的特征值是它作为$H$的自同态时的特征多项式的零点；这个多项式首一且具有整数系数，这就证明了定理的i).当$k=2$且$\chi=1_N$时，可以取奇异同调$H_1(X_0(N),\mathbb Z)$作为$H$.一般情形则需要取``以局部系统为系数''的同调，该局部系统依赖于$k$.\footnote{译注：本段是原文的策略性概述.$H_1(X_0(N),\mathbb Z)$通常是秩为$2g$的同调格，不能直接作为维数$g$的尖点形式空间中的格；其与微分形式及对偶空间的关系将在后文展开.}

第二个想法是先证明$X_0(N)$具有作为$\mathbb Q$上代数曲线的典范模型，然后尝试由此在同调$H$上得到$\operatorname{Gal}(\overline{\mathbb Q}/\mathbb Q)$的作用.这第二步在技术上极为困难，而且并不能很好地实现.我们实际上能做到的，是在以$\mathbb Z/\ell^n\mathbb Z$为系数的同调上定义Galois作用，再取极限，从而在以$\mathbb Z_\ell$为系数的同调$H_{\mathbb Z_\ell}$上定义作用.

至此，我们得到一个有限生成自由$\mathbb Z_\ell$模，其上带有$\mathcal H_0(N)$和$\operatorname{Gal}(\overline{\mathbb Q}/\mathbb Q)$的作用.接下来须证明Galois作用$\rho$在每个$p\nmid\ell N$处非分歧，然后证明在这样的$p$处的Frobenius代换$\sigma_p$满足多项式方程（称为Eichler--Shimura同余关系）
\begin{align*}
1-T(p)\rho(\sigma_p)+pT(p,p)\rho(\sigma_p)^2=0.
\end{align*}
在此基础上，剩下的工作就是证明空间$H_{\mathbb Z_\ell}\otimes_{\mathbb Z_\ell}K_{f,\mathfrak l}[\lambda]$（所有$T_p$的共同特征子空间）在$K_{f,\mathfrak l}$上为二维.这就是所求的表示.\footnote{译注：上述同余关系照录原文.对于所用的同调或上同调约定及算术或几何Frobenius，需要核对是否应对$\rho(\sigma_p)$取逆；不能仅由本段概述将其与此前局部多项式的零点方程直接等同.}

从现在起，假设$k=2$且$\chi=1_N$.因此，我们从$f\in\mathcal S_2(\Gamma_0(N))$出发.在这一情形下，将借助同构
\begin{align*}
\mathcal S_2(\Gamma_0(N))\xrightarrow{\sim}\Omega^1(X_0(N))
:=\{X_0(N)\text{ 上的全纯微分形式 }\omega\},
\end{align*}
尝试完成上述策略.

